\documentclass[11pt]{article}
\usepackage{booktabs}
\usepackage{array}
\usepackage{sectsty}
\usepackage{graphicx}
\usepackage{amsmath,amssymb}
% Margins
\topmargin=-0.45in
\evensidemargin=0in
\oddsidemargin=0in
\textwidth=6.5in
\textheight=9.0in
\headsep=0.25in

\title{A Review of \\Bayesian analysis with conditionally
	identically distributed sequences \\    MR4888176 }
%\author{}
%\date{\today}

\begin{document}
	\maketitle	
	% Optional TOC
	% \tableofcontents
	% \pagebreak





This paper develops a framework for Bayesian-style inference without assuming exchangeability of the observed sequence, instead employing the weaker condition of \emph{conditional identical distribution} (c.i.d.), originally studied by Berti et al. (2004). The central contribution is the demonstration that prior and posterior distributions on a parameter space can still be rigorously defined using c.i.d. sequences, with predictive distributions updated via copula-based martingale constructions.

Let $X_{1:\infty} = (X_1, X_2, \dots)$ be a sequence of real-valued random variables. Standard Bayesian inference relies on exchangeability, which implies the existence of a random probability measure $\mu$ such that, given $\mu$, the observations are i.i.d. (de Finetti's theorem). The authors relax exchangeability to the c.i.d. property, defined by 
\[
\mathbb{E}[g(X_k) \mid X_{1:n}] = \mathbb{E}[g(X_{n+1}) \mid X_{1:n}] \quad \text{a.s.}
\]
for all $k > n \geq 1$ and all bounded measurable $g$. A key consequence is that there still exists a random probability measure $\mu$ such that 
\[
\mu(B) = \lim_{n\to\infty} \mathbb{P}(X_{n+1} \in B \mid X_{1:n}) \quad \text{a.s.},
\]
playing the role of the directing measure.

A central novelty is the construction of prior and posterior distributions on a parameter space $\Theta$ using the maximum likelihood estimator (MLE) from the infinite sequence. Let $f(x;\theta)$ be a parametric model and define 
\[
\hat{\theta}_n = \arg\max_{\theta \in \Theta} \sum_{i=1}^n \log f(X_i;\theta).
\]
Under appropriate regularity conditions (strict concavity, integrability, tail conditions), Theorem 3.2 proves that $\hat{\theta}_n$ converges almost surely to a random variable $\hat{\theta}_\infty$. The distribution of $\hat{\theta}_\infty$ (starting from $X_1$) defines the prior on $\Theta$. Given observed $X_{1:n}$, the sequence $X_{n+1:\infty}$ is again c.i.d., and the distribution of $\hat{\theta}_\infty^{(n)}$ (computed from $X_{n+1:\infty}$) defines the posterior. This mirrors Doob's (1949) insight that posteriors can be derived from predictive distributions without an explicitly specified prior.

To implement the framework, the authors use copulas to model the joint distribution of successive predictions. Theorem 3.1 provides a characterization: $X_{1:\infty}$ is c.i.d. iff there exist copulas $C_{x_{1:n}}$ such that 
\[
\mathbb{P}(X_{n+1} \le x, X_{n+2} \le y \mid X_{1:n}=x_{1:n}) = C_{x_{1:n}}\big(F(x\mid x_{1:n}), F(y\mid x_{1:n})\big),
\]
where 
\[
F(x\mid x_{1:n}) = \mathbb{P}(X_{n+1} \le x \mid X_{1:n}=x_{1:n}).
\]
A practical iterative scheme is 
\[
F_{m+1}(x) = (1-\alpha_m)F_m(x) + \alpha_m C_\rho(F_m(x), F_m(x_{m+1})),
\]
with $\alpha_m = 1/m$ and $C_\rho$ a Gaussian copula:
\[
C_\rho(u,v) = \Phi\left( \frac{\Phi^{-1}(u) - \rho \Phi^{-1}(v)}{\sqrt{1-\rho^2}} \right),
\]
ensuring the predictive distributions $F_m$ form a martingale converging to $\mu$.

A simpler parametric c.i.d. construction is given by 
\[
x_{n} = \theta_{n-1} + \sigma_{n-1}z_n, \quad
\theta_n = (1-a_n)\theta_{n-1} + a_n x_n, \quad
\sigma_n^2 = \sigma_{n-1}^2(1-a_n^2),
\]
with $z_n \stackrel{\text{i.i.d.}}{\sim} N(0,1)$ and $a_n \in (0,1)$. Then $\theta_\infty = \lim \theta_n$ exists almost surely and defines the prior. For mixture models, only the component from which the observation was taken is updated, preserving the c.i.d. property.

The paper provides almost-sure convergence results for the MLE under c.i.d. sequences, extending classical i.i.d. consistency to a more flexible dependence structure. A subtle point raised by a referee is that predictive distributions converge weakly to $\mu$, but not necessarily in total variation; fast enough decay of $\alpha_n$ (e.g., $\sum \alpha_n < \infty$) ensures total variation convergence. This framework may be particularly useful when exchangeability is too restrictive or computationally demanding.

\end{document}



